% \section{Parâmetros Inerciais e Agregação Dinâmica}
% Após definir a geometria do manipulador robótico em 3D através da ferramenta CAD (\emph{Computer-Aided Design}) Autodesk Fusion 360, foi possível obter dados individuais dos elos, das juntas, da base e do manipulador robótico como um todo, este podendo ser referenciado como (\emph{manipulador robótico global}).
% Esses dados individuais e global são as massas, centro de massa local e o tensor de inércia em relação ao próprio centro de massa e frame local, e todos os dados podem ser conferidos no Apêndice A deste trabalho.

% Como o centro de massa é fornecido no frame local do elo, é necessário converter do centro de massa do elo ($\vec{r}_{i}$) para o centro de massa do frame da base ($\vec{r}_{b}$), o termo ($\vec p_i$) é a posição da origem do frame do elo $i$ em relação à base, e o termo ($\vec R_i$) é a matriz de rotação do frame do elo $i$ em relação à base, portanto, utiliza-se:
% % --- Conversão para o frame da base
% \begin{equation}
%     \vec r_b = \vec p_i + \vec R_i \, \vec r_{i},
% \end{equation}

% O tensor de inércia ($\vec I_i$) é dado em relação ao centro de massa no frame do elo, alinhado com seus eixos principais, portanto é necessário rotacioná-lo para o frame da base, conforme:
% \begin{equation}
%     {}^{0}\vec{I}_{i} = \vec{R}_i \,*\, {}^{(elo)}\vec{I}_{i} \,*\, \vec{R}_i^T,
% \end{equation}

\section{Parâmetros Inerciais e Agregação Dinâmica}
Após a modelagem geométrica do manipulador robótico em 3D utilizando a ferramenta CAD (\emph{Computer-Aided Design}) Autodesk Fusion 360, obtiveram-se os parâmetros inerciais individuais de cada elo, junta e base, bem como os parâmetros do manipulador completo. Esses dados compreendem as massas, a posição do centro de massa local e o tensor de inércia em relação ao respectivo centro de massa e \emph{frame} local. Todos os valores detalhados encontram-se no Apêndice A.

\subsection{Parâmetros inerciais por elo}
Como o centro de massa é fornecido no \emph{frame} local do elo, é necessário convertê-lo para o \emph{frame} da base. Denotando por $\vec{r}_{i}$ o vetor do centro de massa no \emph{frame} local do elo, por $\vec{p}_i$ a posição da origem do \emph{frame} do elo $i$ em relação à base, e por $\vec{R}_i$ a matriz de rotação do \emph{frame} do elo $i$ em relação à base, tem-se:
\begin{equation}
    \vec r_b = \vec p_i + \vec R_i \, \vec r_{i}.
\end{equation}

\subsection{Transformação para o frame da base}
De forma análoga, o tensor de inércia ${}^{(elo)}\mathbf{I}_{i}$ fornecido no centro de massa do elo, alinhado com os eixos principais no \emph{frame} local, deve ser rotacionado para o \emph{frame} da base segundo:
\begin{equation}
    {}^{0}\mathbf{I}_{i} =
    \vec{R}_i \,*\,
    {}^{(\text{elo})}\mathbf{I}_{i}
    \,*\, {}^{\, T} \vec{R}_i.
\end{equation}

% Mesmo após rotacionar o tensor de inércia, ele ainda estará sendo referenciado ao respectivo centro de massa, portanto é necessário reescrevê-lo, utilizando o Teorema de Steiner \eqref{eq:Teorema_Steiner}, com relação a origem da base.
% O primeiro termo $\vec I_o$ representa a inércia do elo $i$ em relação à origem da base,
% ${}^0 \vec I_i$ é a inércia no centro de massa rotacionado para o frame da base,
% $m_i$ é a massa do elo $i$, e por fim,
% $\vec r_i$ é a posição do centro de massa em relação à base.
% % --- Teorema de Steiner
% \begin{equation}
%     \label{eq:Teorema_Steiner}
%     \vec{I}_{O_i} = {}^{(0)} \, \vec{I}_{C_i} +
%     m_i \big( \| \vec{r}_i \|^2 \, \vec{I}_3 -
%     \vec{r}_i \,\, \vec{r}_i^{\,T} \big)
% \end{equation}

% Realiza-se a somatória de todas as massas dos elos.
% % --- Agregação total
% \begin{equation}
%     M_{Elos} = \sum_i m_i,
% \end{equation}

% Através de uma média ponderada das posições dos centros de massa de cada elo pode-se obter o centro de massa total do sistema ($CoM_{Global}$)
% \begin{equation}
%     \vec{r}_{Global} = \frac{1}{M_{Elos}} \sum_i m_i \, \vec{r}_i
% \end{equation}

\subsection{Aplicação do Teorema de Steiner}
Mesmo após a rotação, o tensor de inércia ainda está referido ao centro de massa do elo. Assim, aplica-se o Teorema de Steiner \eqref{eq:Teorema_Steiner} para referi-lo à origem da base $O$:
\begin{equation}
  \label{eq:Teorema_Steiner}
  {}^{0} \mathbf{I}_{O_i} =
  {}^{0} \mathbf{I}_{C_i} +
  m_i
    \big(
    ({}^{\,\, T} \vec r_i \,\, \vec r_i)
    \mathbf{I}_3 - \vec r_i \, ({}^{T} \vec r_i)
    \big),
\end{equation}

Onde ${}^{0}\mathbf{I}_{O_i}$ é o tensor de inércia do elo $i$ em relação à origem da base;
${}^{0}\mathbf{I}_{C_i}$ é o tensor no centro de massa já rotacionado para a base;
$m_i$ é a massa do elo $i$;
$\vec r_i$ é a posição do seu centro de massa em relação à base;
e $\mathbf{I}_3$ é a matriz identidade $3\times 3$.

%% ===== EXPLICAÇÃO =====
% A matriz identidade serve para construir um tensor de inercia equivalente que leve em conta a mudança de referencia (do CoM para outro ponto)
% I_3 aparece pq estamos tratando grandezas vetoriais 3D
% é a forma matricial de generalizar a formula escalar I = sum mr^2 para o caso tensorial
% sem ela, não teriamos a componente isotropica


% A massa total (considerando todos os corpos rígidos incluídos no modelo) é:
% \begin{equation}
%   M_{\text{total}} \;=\; \sum_{j} m_j.
% \end{equation}

% O centro de massa global (\emph{CoM} do sistema) é a média ponderada das posições dos centros de massa:
% \begin{equation}
%   \vec r_{\text{Global}} \;=\; \frac{1}{M_{\text{total}}} \sum_{j} m_j \,\vec r_j.
% \end{equation}

% Para encontrar a inércial total do sistema (referindo-se à origem da base) é necessário somar todas as inércias que já foram referidas à origem da base.
% \begin{equation}
%     \vec{I}_O = \sum_i \vec{I}_{O_i},
% \end{equation}

% Para encontrar a inércia total no sistema e que esteja referido ao centro de massa total, é necessário aplicar o Teorema de Steiner novamente:
% \begin{equation}
%     \vec{I}_G = \vec{I}_O -
%     M_{Global} \big(
%     \|\vec{r}_G\|^2 \vec{I}_3 -
%     \vec{r}_G ({}^{\,T}\vec{r}_G) \big),
% \end{equation}

% Para segurança, caso a matriz de inércia não estiver simétrica por conta de erros numéricos, é possível transformá-la em uma matriz de inércia simétrica:
% % --- Simetrização (se necessário)
% \begin{equation}
%     \vec{I} \leftarrow \frac{1}{2} (\vec{I} + \vec{I}^{\, T}),
% \end{equation}

\subsection{Agregação total dos parâmetros do manipulador}
A seguir todos os parâmetros já calculados serão agregados, iniciando pela massa. A massa total (considerando todos os corpos rígidos incluídos no modelo) é:
\begin{equation}
  M_{\text{total}} = \sum_{j} m_j
\end{equation}

O centro de massa global é a média ponderada das posições dos centros de massa:
\begin{equation}
  \vec{r}_{\text{Global}} = \frac{1}{M_{\text{total}}} \sum_{j} m_j \,\vec{r}_j .
\end{equation}

Para encontrar a inércia total do sistema referida à origem da base, soma-se as inércias já referidas à origem:
\begin{equation}
  \mathbf{I}_{O} \;=\; \sum_{i} \mathbf{I}_{O_i} .
\end{equation}

Para obter a inércia total referida ao centro de massa global, aplica-se novamente o Teorema de Steiner, dessa vez transportando a origem $O$ para o centro de massa global:
\begin{equation}
  \mathbf{I}_{G} \;=\; \mathbf{I}_{O} \;-\;
  M_{\text{total}} \Big( \|\vec{r}_{\text{Global}}\|^2 \,\mathbf{I}_3 \;-\;
  \vec{r}_{\text{Global}} \, \, {}^{\,T} \vec{r}_{\text{Global}} \Big) .
\end{equation}

\subsection{Verificações numéricas e consistência}
Uma verificação adicional pode ser feita com o objetivo de verificar a segurança numérica, caso a matriz de inércia apresente assimetria por efeitos de arredondamento, pode-se deixá-la simétrica:
\begin{equation}
  \mathbf{I} \;\leftarrow\; \tfrac{1}{2}\,\big( \mathbf{I} + \mathbf{I}^{\,T} \big).
\end{equation}
